\chapter{Deep Learning}
\label{chap:deeplearning}

Every model in \Cref{chap:classicalml} learns from 320 hand-crafted features. This chapter asks
whether a model that instead learns directly from the raw, filtered signal can do better ---
trained on a separate GPU machine, since the classical pipeline's CPU-only environment has no
PyTorch installed, and scored back here under the same event-level rules as everything else.

\section{Two architectures}

Two compact convolutional architectures were trained, both standard, well-established choices
for EEG decoding rather than custom designs:

\begin{itemize}[leftmargin=1.4em]
  \item \textbf{ShallowConvNet}: a temporal convolution followed by a spatial convolution and
    pooling --- a simple, well-tested baseline architecture for EEG classification.
  \item \textbf{EEGNet}: a compact network built from depthwise and separable convolutions,
    designed specifically to be parameter-efficient on small EEG datasets.
\end{itemize}

Both are deliberately small, with only a few thousand parameters each. With 93 seizures in
total, a larger network would risk simply memorizing the training recordings rather than
learning a general seizure pattern.

\section{Input and training recipe}

Each input is 4 seconds of all 19 filtered EEG channels (1{,}024 samples at 256~Hz), centered on
one of the official 2-second windows, so every CNN prediction lines up with exactly the same
window the classical models were scored on. Each recording is normalized per channel using a
robust statistic, the median and the median absolute deviation, clipped at $\pm 20$, rather than
a plain mean and standard deviation, to reduce the influence of large artifacts.

Training uses the same 5 outer \texttt{GroupKFold} folds as the classical models. Within each
outer fold, some training recordings are further held out for early stopping and for choosing
that fold's detection threshold, so the test fold is never touched during training or threshold
selection. The real, imbalanced class proportions are kept; a weighted cross-entropy loss
up-weights the seizure class to counter the imbalance. Training uses AdamW with a cosine
learning-rate schedule, early stopping on validation average precision, and several random seeds
per fold, with the resulting probabilities averaged across seeds.

Optional online augmentation is applied to \emph{both} classes during training, generated fresh
on the fly per training batch rather than pre-built into a fixed extra dataset --- a
deliberately different strategy from the classical pipeline's augmentation in
\Cref{chap:augmentation}, which only ever touches the seizure class. Eight transforms are
available, each independently selectable so it can be tested alone or combined: the original
five (amplitude scaling, small added noise, channel dropout, left-right mirroring, small time
shifts) plus three more added for this project's wider augmentation sweep: time masking
(zeroing a short random span), frequency band-stop (suppressing a random 1--4~Hz band), and
mixup (linearly blending two windows and their labels).

\section{A real bug: the IndexError in \texttt{event\_counts()}}

During development, the CNN training script crashed with
\texttt{IndexError: index 7943 is out of bounds for axis 0 with size 4344}. The cause: the
function responsible for counting detected events per recording, \texttt{event\_counts()}, used
the \emph{labels} returned by a pandas \texttt{groupby(...).groups} dictionary as if they were
array \emph{positions}. This happens to work only when a table's row index is untouched,
running from 0 to its length minus one. It silently broke the moment the function was handed a
validation subset sliced out of a larger table with \texttt{.iloc[]}, because that subset keeps
its original row numbers rather than renumbering from zero --- and it crashed specifically when
a randomly chosen validation recording's original row number happened to exceed the size of the
smaller validation subset. Earlier runs had not hit this because, by luck, the recordings chosen
for validation had not triggered it.

The fix computes the recording labels, interval starts, and interval ends as plain arrays once,
then selects rows with a boolean mask instead of looking anything up by index label --- correct
regardless of how the table was sliced beforehand. This is the kind of bug that is easy to miss
in testing and only shows up on certain random splits, which is why it is called out here rather
than quietly fixed and forgotten.

\section{Results}

The first run tried the five original transforms only as one fixed combination. A later, wider
sweep tested every transform in isolation --- the same ``test each one alone, then combine the
winners'' approach used on the classical side (\Cref{chap:augmentation}) --- to see whether any
single transform could do better than the original combo, and whether the three new transforms
(time masking, band-stop, mixup) helped at all.

\begin{table}[H]
  \centering
  \small
  \caption{Every CNN run, ShallowConvNet unless noted, nested threshold (best of plain and
    smoothed+hysteresis). FA = false alarms.}
  \label{tab:cnn-results}
  \begin{tabular}{lrrr}
    \toprule
    Model & Event F1 & Found/93 & FA \\
    \midrule
    SVM, channel dropout (classical, new best) & \textbf{0.901} & 88 & 18 \\
    SVM, unbalanced (classical, original baseline) & 0.894 & 86 & 16 \\
    \midrule
    \textbf{Amplitude alone} & \textbf{0.817} & 77 & 19 \\
    Original five, combined (``augmented'') & 0.815 & 77 & 21 \\
    EEGNet, original five combined & 0.811 & 85 & 32 \\
    Mixup & 0.802 & 78 & 24 \\
    Mirror alone & 0.798 & 79 & 27 \\
    Time masking alone & 0.790 & 77 & 25 \\
    Time shift alone & 0.788 & 77 & 26 \\
    Plain (no augmentation) & 0.781 & 76 & 26 \\
    Noise alone & 0.778 & 76 & 27 \\
    Noise + time masking (combined) & 0.774 & 78 & 32 \\
    Band-stop alone & 0.774 & 76 & 28 \\
    Channel dropout alone & 0.767 & 76 & 30 \\
    \bottomrule
  \end{tabular}
\end{table}

None of the 12 CNN runs beats the classical models; the best of them, amplitude scaling alone,
falls 0.084 short of the new SVM best (0.901) and 0.077 short of the original SVM baseline
(0.894). This matches the expectation stated from the start of this project: 93 seizures spread
across 21 patients is too little data for a raw-signal convolutional network to out-learn
hand-designed features. It is recorded here as a real, useful finding, not a failure.

The isolated-transform sweep itself is a smaller, closer race: amplitude scaling alone
(0.817) \emph{barely} edges out the original five-transform combination (0.815) --- a
difference of 0.002, well within normal run-to-run noise, so this should be read as a tie, not a
win for isolating transforms. Unlike the classical pipeline, where channel dropout was a clear,
substantial winner once isolated, no single CNN transform stands out from the pack: most scores
cluster between 0.77 and 0.82 regardless of which transform is used, and combining noise with
time masking (0.774) did not beat either alone. Channel dropout, the classical pipeline's best
method, is actually the \emph{weakest} transform for the CNN (0.767) --- the opposite of what
happened with the feature-based models, a reminder that a transform's effect depends on how a
model is trained, not just on the transform itself.

Smoothing and hysteresis post-processing (described fully in \Cref{chap:evaluation}) was applied
to every CNN run; it helped most of them (plain: $+0.007$; amplitude: $+0.006$; mirror:
$+0.016$; and EEGNet's combo substantially, $+0.035$) but slightly hurt the original five-combo
run ($-0.004$) and time shift ($-0.027$). \Cref{tab:cnn-results} reports each model at whichever
of the two did better for it, which is also how the classical-vs-CNN comparison is framed
throughout this report.

\begin{summarybox}
Two compact CNNs, ShallowConvNet and EEGNet, were trained directly on 4-second raw EEG windows
on a separate GPU machine, under the same recording-level evaluation as the classical models. A
real bug --- an index-label-vs-position mix-up in the event-counting code --- caused a crash on
certain random splits and was fixed once understood. A full sweep of 8 augmentation transforms,
isolated and combined, found no single transform clearly better than the original five-transform
combo (amplitude alone edges it out by 0.002, within noise), and channel dropout --- the
classical pipeline's best method --- is actually the CNN's weakest. None of the 12 CNN runs
(best: 0.817) beats the classical SVM (0.894, or the new best, 0.901), the expected, documented
outcome for a raw-signal model trained on only 93 seizures.
\end{summarybox}
